%% toc: P15 166 166.7
\section{Maximum-entropy existence and reference dependence}
\label{sec:annex-p15-w02b-p15-044-maxent-existence}

Notation is that of \cref{sec:p15-w02b-p15-044-notation}. Every member
\(Q_\lambda\) has a strictly positive \(\nu\)-density, so part~1 of
\cref{thm:p15-w02b-p15-043-gibbs} applies with \(Q=Q_\lambda\).

\begin{annexitem}{annex:p15-w02b-p15-044-attainment}{thm:p15-w02b-p15-044-attainment}{Attainment and interior dual minimizers}
\begin{proof}
\emph{Part 1.} By part~3 of \cref{thm:p15-w02b-p15-043-gibbs},
\(\Psi_c(\lambda)\geq v_\nu(c)=\Psi_c(\lambda^\circ)\) for every
\(\lambda\in\Lambda_\nu\), so \(\lambda^\circ\) minimizes \(\Psi_c\). By
part~2 of that theorem, every \(P\in\mathcal F_\nu(c)\) has defined entropy
and
\[
 H_\nu(P)=\Psi_c(\lambda^\circ)-D(P\Vert Q_{\lambda^\circ}) .
\]
If \(P^*\) is a maximizer, then \(H_\nu(P^*)=v_\nu(c)=\Psi_c(\lambda^\circ)\),
which is finite, so \(D(P^*\Vert Q_{\lambda^\circ})=0\). Part~1 of
\cref{thm:p15-w02b-p15-043-gibbs} gives \(P^*=Q_{\lambda^\circ}\), which is
therefore feasible. Conversely, if \(Q_{\lambda^\circ}\in\mathcal F_\nu(c)\),
then \cref{thm:p15-w02b-p15-043-optimality} makes it the unique maximizer.

\emph{Part 2.} By \cref{thm:p15-w02b-p15-043-optimality},
\(v_\nu(c)=\lambda_0+\lambda\cdot c=\ln Z_\nu(\lambda)+\lambda\cdot c=\Psi_c(\lambda)\).
By part~3 of \cref{thm:p15-w02b-p15-043-gibbs},
\(\Psi_c(\lambda')\geq v_\nu(c)\) for every \(\lambda'\in\Lambda_\nu\), so the
minimum of \(\Psi_c\) is attained at \(\lambda\).

\emph{Part 3.} Let \(\lambda^*\) be interior to \(\Lambda_\nu\) and minimize
\(\Psi_c\) there. Fix \(k\in\{1,\dots,m\}\) and let \(e_k\) be the \(k\)-th
unit vector. Choose \(\varepsilon>0\) with \(\lambda^*+te_k\in\Lambda_\nu\)
for \(|t|<\varepsilon\). By part~2 of \cref{thm:p15-w02b-p15-043-partition},
\(f_k\) is \(Q_{\lambda^*}\)-integrable and \(t\mapsto\ln Z_\nu(\lambda^*+te_k)\)
is differentiable at \(t=0\) with derivative \(-E_{Q_{\lambda^*}}[f_k]\). Hence
\[
 \psi(t)=\Psi_c(\lambda^*+te_k)=\ln Z_\nu(\lambda^*+te_k)+\lambda^*\cdot c+tc_k
\]
is differentiable at \(0\) with \(\psi'(0)=c_k-E_{Q_{\lambda^*}}[f_k]\), and
\(\psi\) has a minimum at \(0\). For \(0<t<\varepsilon\),
\((\psi(t)-\psi(0))/t\geq0\), and for \(-\varepsilon<t<0\) the same quotient
is \(\leq0\); letting \(t\to0\) from each side gives \(\psi'(0)\geq0\) and
\(\psi'(0)\leq0\), so \(\psi'(0)=0\). Thus \(E_{Q_{\lambda^*}}[f_k]=c_k\) for
every \(k\), and \(Q_{\lambda^*}\in\mathcal F_\nu(c)\). By
\cref{thm:p15-w02b-p15-043-optimality} it is the unique maximizer, and by
part~2, \(v_\nu(c)=\Psi_c(\lambda^*)\). No step used regularity of
\(\Lambda_\nu\) or minimality of \(f\).
\end{proof}
\end{annexitem}

\begin{annexitem}{annex:p15-w02b-p15-044-reference}{thm:p15-w02b-p15-044-reference}{Reference change: a fixed-multiplier statistic and an affine test}
\begin{proof}
Since \(\rho>0\), \(\nu'(B)=\int_B\rho\,d\nu=0\) exactly when \(\nu(B)=0\).
Hence \(P\ll\nu\) exactly when \(P\ll\nu'\), and the remaining feasibility
conditions do not involve the reference measure; the feasible sets
coincide. For measurable \(h\geq0\), \(\int h\,d\nu'=\int h\rho\,d\nu\): for
indicators this is the definition of \(\nu'\), for simple functions it
follows by linearity, and for general \(h\) by monotone convergence.

\emph{Part 1.} With \(h=e^{-\lambda\cdot f}\),
\(Z'(\lambda)=\int e^{-\lambda\cdot f}\rho\,d\nu=\int e^{-\lambda\cdot f-(-1)\ln\rho}\,d\nu\),
which is the partition function of \((\nu,(f,\ln\rho))\) at the multipliers
\((\lambda,-1)\). So the two are finite together. When they are, the
\(\nu'\)-density \(e^{-\lambda\cdot f}/Z'(\lambda)\) of \(Q'_\lambda\)
corresponds, by the displayed rule, to the \(\nu\)-density
\(\rho e^{-\lambda\cdot f}/Z'(\lambda)=e^{-\lambda\cdot f-(-1)\ln\rho}/Z'(\lambda)\),
the density of the family member of \((\nu,(f,\ln\rho))\) with multipliers
\((\lambda,-1)\).

\emph{Part 2.} Let \(\ln\rho=b+a\cdot f\) \(\nu\)-almost everywhere. Then
\[
 Z'(\lambda)=\int e^{-\lambda\cdot f+b+a\cdot f}\,d\nu=e^bZ_\nu(\lambda-a),
\]
so \(\lambda\in\Lambda'\) exactly when \(\lambda-a\in\Lambda_\nu\), and
\(\Psi'_c(\lambda)=b+\ln Z_\nu(\lambda-a)+(\lambda-a)\cdot c+a\cdot c
=\Psi_c(\lambda-a)+b+a\cdot c\). By part~1, \(Q'_\lambda\) has
\(\nu\)-density
\(e^{b+a\cdot f-\lambda\cdot f}/(e^bZ_\nu(\lambda-a))=e^{-(\lambda-a)\cdot f}/Z_\nu(\lambda-a)\),
which is the density of \(Q_{\lambda-a}\).

Let \(P\in\mathcal F_\nu(c)\) with \(\nu\)-density \(p\). Since \(P\ll\nu\)
and each \(f_k\) is \(P\)-integrable, \(\ln\rho=b+a\cdot f\) holds
\(P\)-almost everywhere and
\(\int p\,|\ln\rho|\,d\nu=\int|b+a\cdot f|\,dP<\infty\), with
\(\int p\ln\rho\,d\nu=b+a\cdot c\). If \(H_\nu(P)\) is defined,
\cref{thm:p15-w02b-p15-042-reference-change}, applied with reference
\(\nu\) and density \(\rho\), shows that \(H_{\nu'}(P)\) is defined and
equals \(H_\nu(P)+b+a\cdot c\). Conversely \(\nu=\rho^{-1}\nu'\), the
\(\nu'\)-density of \(P\) is \(p/\rho\), and
\(\int(p/\rho)|\ln\rho^{-1}|\,d\nu'=\int|b+a\cdot f|\,dP<\infty\); so the same
theorem, applied with reference \(\nu'\) and density \(\rho^{-1}\), shows that
\(H_\nu(P)\) is defined whenever \(H_{\nu'}(P)\) is. On \(\mathcal F_\nu(c)\)
the two entropies thus have the same domain and differ by the constant
\(b+a\cdot c\), so the two problems have the same maximizers. Their family
members correspond by \(Q'_\lambda=Q_{\lambda-a}\), so the multipliers
differ by \(a\).

\emph{Part 3.} Suppose \(Q'_\lambda=Q_\mu\) for some \(\lambda\in\Lambda'\)
and \(\mu\in\Lambda_\nu\). Their \(\nu\)-densities then agree \(\nu\)-almost
everywhere, by part~1 and uniqueness of densities:
\(\rho e^{-\lambda\cdot f}/Z'(\lambda)=e^{-\mu\cdot f}/Z_\nu(\mu)\). Taking
logarithms,
\(\ln\rho=(\lambda-\mu)\cdot f+\ln Z'(\lambda)-\ln Z_\nu(\mu)\)
\(\nu\)-almost everywhere, which is of the form \(b+a\cdot f\), contrary to
the hypothesis. So no member of one family equals a member of the other.
If each problem has a maximizer in its own family, these maximizers are
members of different families and therefore differ.
\end{proof}
\end{annexitem}

\begin{annexitem}{annex:p15-w02b-p15-044-cubes}{thm:p15-w02b-p15-044-cubes}{The cube example}
\begin{proof}
Lebesgue measure on \((0,1]\) and each \(\nu_s\) have the same null sets,
because the density \((s+1)\ell^s\) is positive on \((0,1]\).

\emph{Part 1.} Regard \(\nu_s\) and the image \(\mu_s\) of Lebesgue measure
on \((0,1]\) under \(u\mapsto u^{1/(s+1)}\) as Borel probability measures on
\(\mathbb R\) giving mass \(0\) outside \((0,1]\). For \(0<a\leq1\),
\(\mu_s((-\infty,a])\) is the length of
\(\{u\in(0,1]:u^{1/(s+1)}\leq a\}=(0,a^{s+1}]\), that is \(a^{s+1}\), and
\(\nu_s((-\infty,a])=\int_0^a(s+1)\ell^s\,d\ell=a^{s+1}\). Both give \(0\)
for \(a\leq0\) and \(1\) for \(a\geq1\). The sets \((-\infty,a]\),
\(a\in\mathbb R\), form a pi-system that generates the Borel sets
(\cref{thm:p15-borel-generators}), so \cref{thm:p15-measure-uniqueness}
gives \(\mu_s=\nu_s\).

\emph{Part 2.} The total mass of \(\nu_s\) is \(1\). Hence
\cref{thm:p15-w02b-p15-043-no-constraint} shows that the unique maximizer
is \(\nu_s\) divided by its total mass, that is, \(\nu_s\) itself. By the computation in part~1,
\(\nu_s((0,1/2])=2^{-(s+1)}\).

\emph{Part 3.} For every \(\lambda\),
\(Z_{\nu_s}(\lambda)=\int e^{-\lambda\ell}\,d\nu_s\leq e^{|\lambda|}\), so
\(\Lambda_{\nu_s}=\mathbb R\), which is open and nonempty; the family is
regular. The statistic \(f(\ell)=\ell\) is minimal: if \(u\ell\) were
\(\nu_s\)-almost everywhere constant with \(u\neq0\), then \(\nu_s\) would
be concentrated on one point, but single points are \(\nu_s\)-null and
\(\nu_s((0,1])=1\). The uniform laws on \((0,1/2]\) and on \((1/2,1]\) are
absolutely continuous with respect to \(\nu_s\) and have means \(1/4\) and
\(3/4\); their mixtures with weights \(1-t\) and \(t\) have means
\(1/4+t/2\). So \([1/4,3/4]\subseteq\mathcal M_{\nu_s}\), and \(1/2\) is an
interior point. Part~1 of \cref{thm:p15-w02b-p15-043-existence} gives
exactly one \(\lambda_s\in\mathbb R\) with \(E_{Q_{\lambda_s}}[\ell]=1/2\),
and \(Q_{\lambda_s}\) is the unique maximizer. Its Lebesgue density is
\((s+1)\ell^se^{-\lambda_s\ell}/Z_{\nu_s}(\lambda_s)\).

Positive constant factors cancel in the mean. Put
\[
 N_s(\lambda)=\int_0^1\ell^{s+1}e^{-\lambda\ell}\,d\ell,\qquad
 D_s(\lambda)=\int_0^1\ell^se^{-\lambda\ell}\,d\ell>0;
\]
the mean is \(m_s(\lambda)=N_s(\lambda)/D_s(\lambda)\). So \(\lambda_s\) is
the unique solution of the displayed equation. If \(\lambda_n\to\lambda\)
with \(|\lambda_n|\leq M\), the integrands of \(N_s\) and \(D_s\) are
bounded by \(e^M\) on \((0,1]\), and dominated convergence makes \(N_s\),
\(D_s\) and \(m_s\) continuous. At \(\lambda=0\),
\(m_s(0)=(1/(s+2))/(1/(s+1))=(s+1)/(s+2)\), which is \(1/2\), \(2/3\) and
\(3/4\). Hence \(\lambda_0=0\).

For \(s=1,2\) we bound \(m_s(10)\). By \cref{thm:p15-gamma-normalization},
the \(\operatorname{gamma}(n+1,10)\) density integrates to \(1\), so
\(\int_0^\infty\ell^ne^{-10\ell}\,d\ell=n!/10^{n+1}\). Differentiation shows
that \(-e^{-10\ell}(\ell/10+1/100)\) and
\(-e^{-10\ell}(\ell^2/10+\ell/50+1/500)\) are antiderivatives of
\(\ell e^{-10\ell}\) and \(\ell^2e^{-10\ell}\), and both tend to \(0\) as
\(\ell\to\infty\). Hence
\[
 \int_1^\infty\ell e^{-10\ell}\,d\ell=0.11\,e^{-10},\qquad
 \int_1^\infty\ell^2e^{-10\ell}\,d\ell=0.122\,e^{-10}.
\]
Since \(e>2.7\) and \(2.7^{10}>20000\), both are less than \(10^{-5}\).
Bounding each numerator by the integral over \((0,\infty)\),
\[
 m_2(10)\leq\frac{6\cdot10^{-4}}{2\cdot10^{-3}-10^{-5}}<\frac12,
 \qquad
 m_1(10)\leq\frac{2\cdot10^{-3}}{10^{-2}-10^{-5}}<\frac12 .
\]
Since \(m_s(10)<1/2<m_s(0)\), the intermediate value theorem gives a root
in \((0,10)\), and by uniqueness it is \(\lambda_s\).

\emph{Part 4.} Relative to one another the references are
\(\nu_1=\rho_{10}\nu_0\), \(\nu_2=\rho_{20}\nu_0\) and
\(\nu_2=\rho_{21}\nu_1\) with \(\rho_{10}=2\ell\), \(\rho_{20}=3\ell^2\)
and \(\rho_{21}=3\ell/2\), so each \(\ln\rho\) has the form
\(\ln k+j\ln\ell\) with \(k>0\) and \(j\in\{1,2\}\). Such a function is not
\(\nu_s\)-almost everywhere equal to \(b+a\ell\). Otherwise the two
continuous functions on \((0,1]\) would agree outside a null set, hence on a
dense set, since every nonempty open subinterval has positive measure, and
so everywhere by continuity. That is impossible, because
\(\ln k+j\ln\ell\to-\infty\) as \(\ell\downarrow0\) while \(b+a\ell\) stays
bounded. Each maximizer of part~3 belongs to its own family, so part~3 of
\cref{thm:p15-w02b-p15-044-reference} shows that they are pairwise distinct.
Finally, the \(\nu_0\)-density of the \(\nu_2\) maximizer is
\(3\ell^2e^{-\lambda_2\ell}/Z_{\nu_2}(\lambda_2)
 =\exp\{-\lambda_2\ell-(-2)\ln\ell\}\big/\bigl(Z_{\nu_2}(\lambda_2)/3\bigr)\).
This is a probability density with respect to \(\nu_0\), so the
denominator equals the \(\nu_0\) partition function of \((\ell,\ln\ell)\) at
\((\lambda_2,-2)\), and the maximizer is that family member, as part~1 of
\cref{thm:p15-w02b-p15-044-reference} also states.

The decimal values in \cref{def:p15-w02b-p15-044-cube-example} were
obtained by solving \(m_s(\lambda)=1/2\) numerically in thirty-digit
arithmetic and integrating the resulting densities over \((0,1/2]\); they
lie in the interval \((0,10)\) established above.
\end{proof}
\end{annexitem}
